\section{From Mathematics to Code: A Three-Tier Derivation}\label{sec:m2c}

The contribution of this report rests on a deliberate three-tier
descent: from the mathematical operator~\(\mathcal{G}\) of
\cref{sec:soundness}, to a deterministic algorithm, to a Linux-kernel
implementation path. Each tier preserves the invariants of the tier
above, by construction. This section presents the descent in a single
narrative arc, ending with a traceability table
(\cref{tab:traceability}) that maps each line of the proof to the
algorithmic step and the kernel construct that realizes it. The aim
is auditability: a reviewer should be able to follow the chain
\emph{theorem~\(\rightarrow\)~pseudocode~\(\rightarrow\)~kernel call} line
by line.

\subsection{Tier~1: The Recovery Operator (Mathematics)}\label{sec:m2c-tier1}

\begin{definition}[Recovery operator]\label{def:G}
The BEAMFS recovery operator~\(\mathcal{G} : \mathcal{S} \to
\mathcal{S} \cup \{\bot\}\) is the partial function defined by the
following procedure on input~\(s = (s_V, s_P)\):
\begin{enumerate}
\item For every region~\(R\) of~\(\Lambda\), evaluate the redundancy
  syndrome on~\(s|_R\). If all syndromes are zero, return~\(s\)
  (idempotence by~I4).
\item Otherwise, for the first region~\(R^{*}\) (in canonical region
  order) whose syndrome is non-zero, attempt RS correction
  on~\(s|_{R^{*}}\) using the parity bytes adjacent to that region.
  Let~\(s|_{R^{*}}'\) denote the corrected bytes.
\item Re-evaluate CRC on~\(s|_{R^{*}}'\). If CRC fails, return~\(\bot\)
  (corruption exceeded RS correction capacity, by~I1).
\item If~\(R^{*}\) is a region with both layers, re-establish coherence
  by overwriting the lagging copy with the corrected leading copy.
\item Verify all sentinels in~\(R^{*}\) are at canonical values. If
  any sentinel is corrupted in a way RS missed, return~\(\bot\).
\item Append a journal entry to the RAF journal documenting the
  correction (block number, timestamp, symbol count, entropy
  estimate where applicable).
\item Recurse: \(\mathcal{G}\bigl(\text{state with } s|_{R^{*}}
  \text{ replaced by } s|_{R^{*}}'\bigr)\).
\end{enumerate}
\end{definition}

\begin{proposition}[Idempotence on consistent inputs]\label{prop:idem}
For every~\(s \models\),~\(\mathcal{G}(s) = s\).
\end{proposition}
\begin{proof}[Proof sketch]
By \cref{def:consistent}, all syndromes are zero on a consistent
state. Step~1 of \cref{def:G} returns~\(s\) without modification.
\end{proof}

\begin{lemma}[Descent under perturbation]\label{lem:descent}
Let~\(N(s)\) denote the number of regions in~\(\Lambda\) whose
syndrome is non-zero on~\(s\). Then for every~\(s\),
either~\(\mathcal{G}(s) = \bot\), or
\(N\bigl(\mathcal{G}(s)\bigr) \le N(s) - 1.\)
\end{lemma}
\begin{proof}[Proof sketch]
Inspect the cases of \cref{def:G}. Step~1 with~\(N(s) = 0\) yields
the same state; if~\(N(s) \ge 1\), the loop body either
returns~\(\bot\) (cases~3, 5) or replaces a corrupt region by a
CRC-verified, sentinel-verified consistent region, decreasing~\(N\)
by exactly one before the recursive call.
\end{proof}

\begin{proposition}[Fail-closed absorbing]\label{prop:absorb}
\(\mathcal{G}(\bot) = \bot\). If~\(\mathcal{G}(s) = \bot\), then for
every~\(s'\) reachable from~\(s\) by a sequence of
\(\mathcal{G}\)-applications, \(\mathcal{G}(s') = \bot\) (invariant~I5).
\end{proposition}
\begin{proof}[Proof sketch]
Direct from \cref{def:G}: \(\bot\) has no syndrome to evaluate; the
procedure terminates at step~1 by convention, returning~\(\bot\).
I5 then propagates by induction on the application chain.
\end{proof}

These three results, combined with \cref{thm:soundness} of
\cref{sec:soundness}, fully characterize~\(\mathcal{G}\) as a
mathematical object. They are the contract that the next two tiers
must respect.

\subsection{Tier~2: The Algorithm}\label{sec:m2c-tier2}

We translate~\(\mathcal{G}\) into a deterministic algorithm in
C-style pseudocode (kernel-conformant: ISO C with kernel
extensions, no C++ idioms). Each numbered step of \cref{def:G}
corresponds to a labelled comment in the algorithm.
\Cref{tab:traceability} below makes the correspondence explicit.

\begin{lstlisting}[language=C,frame=single,
  basicstyle=\ttfamily\footnotesize]
/*
 * BEAMFS recovery operator G, single application.
 * Returns 0 on success (state consistent or unchanged),
 * 1 if a region was corrected, -EBADMSG on fail-closed.
 *
 * Caller iterates until return value is 0 or -EBADMSG.
 */
int beamfs_g_step(struct beamfs_state *s)
{
    struct region *R;
    u32 syndrome;
    int rc;

    list_for_each_entry(R, &s->lambda, list) {
        syndrome = rs_syndrome(s, R);
        if (syndrome == 0)
            continue;                /* I4: skip consistent */

        /* Step 2: RS correction */
        rc = rs_correct_region(s, R);
        if (rc < 0)
            return -EBADMSG;          /* uncorrectable */

        /* Step 3: re-verify CRC */
        if (crc32_region(s, R) != stored_crc(s, R))
            return -EBADMSG;          /* RS produced wrong codeword */

        /* Step 4: re-establish coherence */
        if (region_is_dual_layer(R))
            propagate_corrected(s, R);

        /* Step 5: sentinel verification */
        if (!sentinels_canonical(s, R))
            return -EBADMSG;

        /* Step 6: journal */
        raf_log_event(s, R, rc /* symbol_count */);

        return 1;                     /* one region corrected */
    }
    return 0;                          /* fixpoint: all consistent */
}

/*
 * Driver loop: at most |Lambda| iterations by Lemma 5.3.
 */
int beamfs_recover(struct beamfs_state *s)
{
    int iter = 0, rc;

    while ((rc = beamfs_g_step(s)) == 1) {
        if (++iter > LAMBDA_MAX)
            return -EBADMSG;          /* belt-and-braces */
    }
    return rc;                         /* 0 or -EBADMSG */
}
\end{lstlisting}

\paragraph{Complexity.}
Each call to \texttt{beamfs\_g\_step} performs at most one RS decode
on the affected region plus constant-time syndrome checks elsewhere.
The dominant cost is the RS decode, \(O(n^{2})\) for a
Berlekamp--Massey decoder over GF(\(2^{8}\)) with \(n=255\),
approximately \(65{,}000\) GF operations per subblock. The driver
loop runs at most \(|\Lambda|\) iterations
(\cref{lem:descent}), so total worst-case complexity is
\(O(|\Lambda| \cdot n^{2})\). For the FTRFS layout, counting the
five named regions of \cref{def:layout} together with their
constituent RS subblocks yields~\(|\Lambda| \approx 16\); with
\(n^{2} = 65{,}025\) this gives approximately one million GF
operations per worst-case full recovery---a sub-millisecond cost
on contemporary CPUs. The bound is robust to a factor of a few
in the subblock count.

\subsection{Tier~3: The Kernel Implementation}\label{sec:m2c-tier3}

The reference implementation of BEAMFS targets the Linux kernel
through the standard VFS infrastructure
(\texttt{file\_system\_type}, \texttt{iomap}). The recovery
operator runs kernel-side, invoked at mount
(\texttt{ftrfs\_fill\_super}) and on every metadata-region read
path. Userspace recovery is excluded by~I3: the coherence relation
between volatile and persistent layers cannot be enforced from a
process holding its own page-cache view.

\paragraph{Language choice.}
The reference implementation is in C, compiled by the kernel
toolchain. C is the kernel's lingua franca; upstreaming requires
it. The RS hot path reuses the kernel's audited
\texttt{lib/reed\_solomon}~\cite{kernelrslib}, preserving the
audit-surface argument of FTRFS v1~\cite{desbrieres2026ftrfs}.
A future v2 may explore a Rust port of the driver loop (not the
GF arithmetic), following the model
of~\cite{sigurbjarnarson2016yggdrasil} and the recent
kernel Rust support.

\paragraph{Differences from FTRFS v1.}
FTRFS v1 implements a proto-version of the operator without the
formal calculus. A BEAMFS-compliant implementation introduces:
\begin{itemize}
\item Explicit invariant table compiled into the module
  (\texttt{lambda[]}), making region iteration order canonical and
  auditable.
\item Mandatory journal entry on every correction event, with the
  format-v4 entropy field replacing FTRFS v3's coarse symbol count.
\item Driver loop explicitly bounded at~\(|\Lambda|\); refusing to
  mount a layout where this bound is violated by the data.
\item Fail-closed transition flipping the superblock to read-only
  mode; subsequent mutations rejected until offline repair
  (\texttt{fsck.beamfs}, planned).
\end{itemize}

\subsection{Traceability: Math~\texorpdfstring{$\rightarrow$}{to}~Algorithm~\texorpdfstring{$\rightarrow$}{to}~Code}\label{sec:m2c-trace}

\Cref{tab:traceability} expresses the descent in tabular form.
Each row corresponds to one mathematical commitment, the
algorithmic line that realizes it, and the kernel construct that
will instantiate it. The table is the auditable contract: a kernel
patch series proposing BEAMFS must exhibit each kernel construct
of the third column, and a reviewer must be able to verify that
each construct enforces the property of the first.

\begin{table*}[t]
\centering\small
\caption{Traceability: each mathematical commitment in the operator
definition (Tier~1) maps to an algorithmic step (Tier~2) and a
planned kernel construct (Tier~3). The phase indicates when the
construct is checked: \emph{Static} = compile-time
(\texttt{BUILD\_BUG\_ON} or \texttt{static\_assert});
\emph{Mount} = once at \texttt{fill\_super}; \emph{Runtime} = on
each metadata access.}
\label{tab:traceability}
\begin{tabularx}{\textwidth}{l X X X l}
\toprule
\textbf{Math} & \textbf{Commitment (Tier~1)} & \textbf{Algorithm (Tier~2)} & \textbf{Kernel construct (Tier~3)} & \textbf{Phase} \\
\midrule
I1 (Total redundancy) &
each region's RS capacity covers worst-case MBU width &
\texttt{rs\_correct\_region(s, R)} returns \texttt{<0} when \(>t\) errors &
\texttt{rs\_subblock\_capacity[]} table; \texttt{BUILD\_BUG\_ON} per region &
Static + Mount \\
\addlinespace[0.3em]
I2 (Sentinel separation) &
canonical-value bytes detectable in syndrome alone &
\texttt{sentinels\_canonical(s, R)} verification &
zero-pad fields in on-disk structs; \texttt{BUILD\_BUG\_ON(offsetof(\dots))} &
Static \\
\addlinespace[0.3em]
I3 (Coherence checkpoint) &
volatile/persistent serialization total, deterministic, bijective &
\texttt{propagate\_corrected(s, R)} on dual-layer regions &
\texttt{ftrfs\_dirty\_super()} family; single-source-of-truth canonical serializer &
Runtime \\
\addlinespace[0.3em]
I4 (Idempotence) &
\(\mathcal{G}(s) = s\) for every \(s \models\) &
\texttt{if (syndrome == 0) continue;} early skip in inner loop &
\texttt{rs\_syndrome() == 0} early return &
Runtime \\
\addlinespace[0.3em]
I5 (Fail-closed monotonicity) &
\(\mathcal{G}(\bot) = \bot\); no recovery from \(\bot\) &
\texttt{return -EBADMSG} (steps~3, 5) absorbing &
\texttt{set\_sb\_rdonly()} + journal mark; subsequent mutations rejected &
Runtime \\
\addlinespace[0.3em]
\Cref{def:G}, step~1 &
syndrome evaluation per region &
\texttt{u32 syndrome = rs\_syndrome(s, R);} &
\texttt{ftrfs\_rs\_syndrome\_region()} in \texttt{super.c} &
Runtime \\
\addlinespace[0.3em]
\Cref{def:G}, step~2 &
RS correction on first non-zero region &
\texttt{rs\_correct\_region(s, R)} &
\texttt{decode\_rs8()} from \texttt{lib/reed\_solomon}~\cite{kernelrslib} &
Runtime \\
\addlinespace[0.3em]
\Cref{def:G}, step~3 &
CRC re-verification after RS correction &
\texttt{if (crc32\_region(s, R) != stored\_crc(s, R))} &
\texttt{crc32\_le()} from \texttt{lib/crc32.c} &
Runtime \\
\addlinespace[0.3em]
\Cref{def:G}, step~6 &
journal entry on every correction event &
\texttt{raf\_log\_event(s, R, rc)} &
RAF journal v4 entry append (\texttt{ftrfsd}) &
Runtime \\
\addlinespace[0.3em]
\Cref{lem:descent} &
\(N(\mathcal{G}(s)) \le N(s) - 1\) per non-fail-closed step &
each \texttt{return 1} reduces non-zero region count by one &
loop counter bounded at \texttt{LAMBDA\_MAX} &
Runtime \\
\bottomrule
\end{tabularx}
\end{table*}

\subsection{Why a Calculus, Not a Procedure}\label{sec:m2c-calculus}

A reviewer may legitimately ask: why frame the work as a
\emph{calculus} rather than as a recovery procedure? The
distinction matters for the descent presented above. A procedure
is verifiable case by case but composes unpredictably---each
combination of cases requires re-verification. A calculus is a
small set of operators with stated algebraic properties
(idempotence by~\cref{prop:idem}, descent
by~\cref{lem:descent}, fail-closed absorption
by~\cref{prop:absorb}); its composition is governed by those
properties, once and for all. The calculus is what allows
\cref{thm:soundness} to be stated globally, rather than re-proved
per call site. It is also what allows future filesystems beyond
BEAMFS, possibly with different layouts, to inherit the soundness
result by exhibiting an instance of the same calculus---without
re-proving the descent.

\subsection{The Soundness--Implementation Gap}\label{sec:m2c-gap}

The three tiers above preserve the invariants by construction. The
remaining gap is between Tier~3 and the binary actually executed by
the kernel: the standard verified-compilation gap. Closing it
requires a verified compilation chain in the seL4
sense~\cite{klein2009sel4} or a verified
filesystem in the FSCQ sense~\cite{chen2015fscq}, neither of
which is in scope for v1. v3 of this report will discuss
mechanization of \cref{thm:soundness} in a proof assistant
(Coq, Isabelle/HOL, or Lean) and the corresponding implications
for the implementation gap.
